% Part: methods
% Chapter: induction
% Section: inductive-definitions

\documentclass[../../../include/open-logic-section]{subfiles}

\begin{document}

\olfileid{mth}{ind}{idf}

\olsection{استقرايي تعريفونه}

په منطق کښې ډېر ځله د څيزونو ډولونه \emph{په استقرايي ډول}
تعريفوو؛ يعنې د دې دپاره چې څه د تعريفېدونکي ډول څيز ګڼل
کېږي، داسې قواعد ټاکو چې بيانوي څنګه د همدغه ډول له پخوانيو
څيزونو څخه نوي څيزونه ترلاسه کړو۔ د بېلګې په توګه، ډېر
ځله د نښو د لړيو ځانګړي ډولونه، لکه د يوې ژبې ترمونه او
!!{formula}s، په استقرا تعريفوو۔ د يوې ساده بېلګې دپاره
داسې تارونه په نظر کښې ونيسئ چې له $\mathrm{a}$، $\mathrm{b}$،
$\mathrm{c}$، $\mathrm{d}$ تورو، له~$\circ$ نښې او له $[$ او
$]$ قوسونو جوړ دي، لکه «$[[\mathrm{c} \circ \mathrm{d}][$»،
«$[\mathrm{a}[]\circ]$»، «$\mathrm{a}$» يا
«$[[\mathrm{a} \circ \mathrm{b}]\circ
\mathrm{d}]$»۔ ښايي احساس وکړئ چې په لومړيو دوو تارونو
کښې څه «ناسم» دي: په لومړي کښې قوسونه هېڅ «برابر» نۀ
دي، او ښايي احساس وکړئ چې «$\circ$» بايد داسې عبارتونه
«سره وتړي» چې پخپله معنا ولري۔ درېيم او څلورم تار ښه
ښکاري: د هر «$[$» دپاره يو تړونکے «$]$» شته، که قوسونه
پکښې وي، او د هر $\circ$ دواړو خواوو ته «ښه» عبارتونه
موندلاے شو چې د يوې جوړې قوسونو دننه دي۔

غواړو په کره توګه وټاکو چې څه «ښه‌جوړ ترم» ګڼل کېږي۔
تر ټولو وړاندې، هر تورے په يوازې ځان ښه‌جوړ دے۔ هر څه
چې يوازې يو تورے نۀ وي، بايد د «$[t \circ s]$» په بڼه
وي، چې $s$ او $t$ پخپله ښه‌جوړ وي۔ برعکس، که $t$ او
$s$ ښه‌جوړ وي، نو د دواړو تر منځ د يو $\circ$ په
اېښودلو او د يوې جوړې قوسونو دننه ئې په راوستلو يو
نوی ښه‌جوړ ترم جوړولاے شو۔ له دې عملونو څخه د ښه‌جوړو
ترمونو د سټ د \emph{تعريفولو} دپاره کار اخيستلاے شو۔
دا يو \emph{استقرايي تعريف} دے۔

\begin{defn}[ښه‌جوړ ترمونه]
د \emph{ښه‌جوړو ترمونو} سټ په استقرايي ډول داسې تعريفېږي:
  \begin{enumerate}
  \item هر $\mathrm{a}$، $\mathrm{b}$، $\mathrm{c}$،
    $\mathrm{d}$ تورے يو ښه‌جوړ ترم دے۔
  \item که $s_1$ او $s_2$ ښه‌جوړ ترمونه وي، نو
    $[s_1 \circ s_2]$ هم ښه‌جوړ ترم دے۔
  \item نور هېڅ څه ښه‌جوړ ترم نۀ دي۔
  \end{enumerate}
\end{defn}

دا تعريف راته وايي چې يو څيز هله او يوازې هله ښه‌جوړ
ترم ګڼل کېږي چې د (1) او~(2) دوو شرطونو له مخې په يو
متناهي شمېر ګامونو کښې جوړېداے شي۔ په لومړي ګام کښې
ټول هغه ښه‌جوړ ترمونه جوړوو چې يوازې تورې دي، يعنې:
\[
\mathrm{a}, \mathrm{b}, \mathrm{c}, \mathrm{d}
\]
په دويم ګام کښې (2) پر جوړو شوو ترمونو کاروو۔ ترلاسه کوو:
\[
[\mathrm{a} \circ \mathrm{a}], [\mathrm{a} \circ \mathrm{b}],
[\mathrm{b} \circ \mathrm{a}], \dots, [\mathrm{d} \circ \mathrm{d}]
\]
دا د دوو تورو د ټولو يوځاے کېدو دپاره دي۔ په درېيم ګام
کښې بيا (2) کاروو، پر هر دوو ښه‌جوړو ترمونو چې تر اوسه
مو جوړ کړي دي۔ نوي ښه‌جوړ ترمونه ترلاسه کوو، لکه
$[\mathrm{a} \circ [\mathrm{a} \circ
    \mathrm{a}]]$، چې $t$ پکښې $\mathrm{a}$ د لومړي ګام
څخه او $s$ پکښې $[\mathrm{a} \circ \mathrm{a}]$ د دويم
ګام څخه دے؛ او $[[\mathrm{b} \circ
    \mathrm{c}] \circ [\mathrm{d} \circ \mathrm{b}]]$، چې
له دوو ترمونو $[\mathrm{b} \circ \mathrm{c}]$ او $[\mathrm{d}
  \circ \mathrm{b}]$ څخه د دويم ګام په بنسټ جوړ شوے دے۔
او همداسې نور۔ (3) شرط دا امکان ردوي چې کوم څيز چې
په دې ډول نۀ وي جوړ شوے، د ښه‌جوړو ترمونو سټ ته ننوځي۔

په ياد ولرئ چې تر اوسه مو دا نۀ دي ثابت کړي چې د نښو
هره لړۍ چې ښه‌جوړه «احساسېږي»، د دې تعريف له مخې هم
ښه‌جوړه ده۔ خو دا بايد ښکاره وي چې هر څه چې جوړولاے
شو، په رښتيا ښه‌جوړ «احساسېږي»: قوسونه برابر دي، او
$\circ$ داسې برخې سره تړي چې پخپله ښه‌جوړې دي۔

د استقرايي تعريفونو مهمه ځانګړنه دا ده چې که د ټولو
ښه‌جوړو ترمونو په هکله څه ثابتول غواړئ، تعريف راته
وايي چې کوم حالتونه بايد په نظر کښې ونيسو۔ د بېلګې
په توګه، که درته وويل شي چې $t$ يو ښه‌جوړ ترم دے، نو
استقرايي تعريف راته وايي چې $t$ څنګه کېداے شي: $t$
يو تورے کېداے شي، يا $[s_1 \circ s_2]$ کېداے شي، د
ښه‌جوړو ترمونو د يوې جوړې $s_1$ او~$s_2$ دپاره۔ د
(3) شرط له امله همدا يوازيني امکانات دي۔

د يو په استقرايي ډول تعريف شوي سټ د ټولو غړو په هکله
د دعوو د ثابتولو پر وخت د استقرا قوي بڼه په ځانګړي
ډول مهمه کېږي۔ د بېلګې په توګه، فرض کړئ غواړو ثابت
کړو چې د اوږدوالي~$n$ په هر ښه‌جوړ ترم کښې د $[$
شمېر~$< n/2$ دے۔ دا د ټولو~$n$ په هکله يوه دعوه
ګڼلاے شو: د هر $n$ دپاره، د $[$ شمېر په هر هغه
ښه‌جوړ ترم کښې چې اوږدوالے ئې~$n$ دے، $< n/2$ دے۔

\begin{prop}
د هر $n$ دپاره، د $[$ شمېر په هغه ښه‌جوړ ترم کښې چې
اوږدوالے ئې~$n$ دے، $< n/2$ دے۔
\end{prop}

\begin{proof}
د قوي استقرا په ذريعه د دې نتيجې د ثابتولو دپاره
بايد وښيو چې دا شرطي دعوه رښتيا ده:
\begin{quote}
که د هر $l < k$ دپاره، د اوږدوالي~$l$ په هر ښه‌جوړ ترم
کښې د نښو شمېر $< l/2$ وي، د $[$ نښو، نو د
اوږدوالي~$k$ په هر ښه‌جوړ ترم کښې ئې شمېر $< k/2$
دے، د $[$ نښو۔
\end{quote}
د دې شرط د ښودلو دپاره فرض کړئ چې مقدم ئې رښتيا
دے؛ يعنې د هر $l<k$ دپاره، د اوږدوالي~$l$ په ښه‌جوړو
ترمونو کښې د نښو شمېر $< l/2$ دے، د $[$ نښو۔ دې فرض ته
استقرايي فرض وايو۔ غواړو وښيو چې همدا خبره د
اوږدوالي~$k$ ښه‌جوړو ترمونو دپاره هم رښتيا ده۔

نو فرض کړئ چې $t$ يو ښه‌جوړ ترم دے چې اوږدوالے
ئې~$k$ دے۔ ځکه چې ښه‌جوړ ترمونه په استقرايي ډول
تعريف شوي دي، دوه حالتونه لرو: (1)~$t$ يوازې يو
تورے دے، يا (2)~$t$، $[s_1 \circ s_2]$ دے، د ځينو
ښه‌جوړو ترمونو $s_1$ او~$s_2$ دپاره۔
\begin{enumerate}
\item $t$ يو تورے دے۔ بيا $k = 1$، او د $[$ شمېر په
$t$ کښې~$0$ دے۔ ځکه چې $0 < 1/2$، دعوه سمه ده۔
\item $t$، $[s_1 \circ s_2]$ دے، د ځينو ښه‌جوړو ترمونو
  $s_1$ او~$s_2$ دپاره۔ $l_1$ دې د~$s_1$ اوږدوالے
  او $l_2$ دې د~$s_2$ اوږدوالے وي۔ بيا اوږدوالے~$k$،
  د $t$ دپاره، $l_1+l_2+3$ دے؛ د $s_1$ او $s_2$
  اوږدوالے او درې نښې $[$، $\circ$، $]$۔ ځکه چې
  $l_1+l_2+3$ تل له $l_1$ څخه زيات دے، $l_1 < k$۔
  همداسې $l_2
  < k$۔ دا معنا لري چې استقرايي فرض پر ترمونو
  $s_1$ او~$s_2$ کارېږي: شمېر~$m_1$، د $[$ دپاره
  په $s_1$ کښې، $< l_1/2$ دے، او شمېر~$m_2$، د
  $[$ دپاره په~$s_2$ کښې، $< l_2/2$ دے۔

د $[$ شمېر په $t$ کښې د $[$ د شمېر په~$s_1$
کښې، د $[$ د شمېر په~$s_2$ کښې، او د~$1$ مجموعه
ده؛ يعنې $m_1 + m_2 + 1$ دے۔ ځکه چې $m_1
  < l_1/2$ او $m_2 < l_2/2$، نو لرو:
  \[
  m_1 + m_2 + 1 < \frac{l_1}{2} + \frac{l_2}{2} + 1 = \frac{l_1+l_2+2}{2} < \frac{l_1+l_2+3}{2} = k/2.
  \]
\end{enumerate}
په هر حالت کښې مو ښودلې چې د $[$ شمېر په $t$ کښې
$< k/2$ دے، د استقرايي فرض په بنسټ۔ په قوي استقرا
مطلوبه دعوه ثابتېږي۔
\end{proof}

\begin{prob}
د ډېرو ښه‌جوړو ترمونو سټ داسې تعريف کړئ:
  \begin{enumerate}
  \item هر $\mathrm{a}$، $\mathrm{b}$، $\mathrm{c}$،
    $\mathrm{d}$ تورے يو ډېر ښه‌جوړ ترم دے۔
  \item که $s$ ډېر ښه‌جوړ ترم وي، نو $[s]$ هم دے۔
  \item که $s_1$ او $s_2$ ډېر ښه‌جوړ ترمونه وي، نو
    $[s_1 \circ s_2]$ هم دے۔
  \item نور هېڅ څه ډېر ښه‌جوړ ترم نۀ دي۔
  \end{enumerate}
وښيئ چې د $[$ شمېر په ډېر ښه‌جوړ ترم~$t$ کښې، چې
اوږدوالے ئې~$n$ دے، $\le n/2 +1$ دے۔
\end{prob}

\end{document}
